\documentclass[11pt]{article}

\usepackage[a4paper,margin=0.92in]{geometry}
\usepackage{amsmath,amssymb,amsthm,mathtools}
\usepackage{booktabs}
\usepackage{graphicx}
\usepackage{xcolor}
\usepackage[colorlinks=true,linkcolor=blue!55!black,citecolor=blue!55!black,urlcolor=blue!55!black]{hyperref}
\usepackage{microtype}

\newtheorem{theorem}{Theorem}
\newtheorem{lemma}{Lemma}
\newtheorem{proposition}{Proposition}
\newtheorem{corollary}{Corollary}
\newtheorem{remark}{Remark}
\newcommand{\R}{\mathbb{R}}
\newcommand{\BV}{BV_{\mathrm{loc}}}
\newcommand{\HH}{\mathcal{H}}
\newcommand{\V}{\mathsf{V}}
\newcommand{\Tmap}{\mathcal{T}}
\newcommand{\Cb}{C_b}
\DeclareMathOperator{\dist}{dist}

\title{Variability from Bounded Compression and Transversal Fixed Points}
\author{Automatic Math Discovery Run \texttt{fa\_banach\_001}\\Model: GPT5.6}
\date{26 August 2026}

\begin{document}
\maketitle

\begin{center}
\fcolorbox{blue!45!black}{blue!4}{\parbox{0.91\textwidth}{
\textbf{Status.} Candidate substantial partial result, likely valid.  The
bounded-compression theorem covers general $BV$ coefficients for almost every
initial point and the interface theorems give explicit fixed-point classes.
The packet still does not replace the Doss transform for a general rough
driver at every prescribed initial point.
}}
\end{center}

\begin{abstract}
Hinz, T\"olle, and Viitasaari ask for assumptions under which the integral
process for a $BV$ coefficient is itself variable.  The main new observation
is measure theoretic: if a family of trajectories has uniformly bounded
compression, then for Lebesgue-almost every initial point its trajectory is
$(s,1)$-variable with respect to every locally $BV$ coefficient, simultaneously
for all $s<1$.  Applied to the source's bi-Lipschitz Doss flow, this removes
all occupation assumptions on an arbitrary H\"older driver for almost every
initial point.  Applied to Ambrosio's regular Lagrangian flow theorem, it gives
a non-Doss result for bounded $BV$ coefficients with bounded column
divergences and Lipschitz drivers.  We also retain two pointwise fixed-point
theorems for a genuine hyperplane jump.  A common normal row and a
bi-Lipschitz driver clock yield Young--Picard solutions for rough drivers; a
positive normal velocity yields a direct Schauder solution for absolutely
continuous drivers.  Both have the sharp conclusion
$X\in\mathsf V(\sigma,s,p)$ for $sp<1$.  An exact elliptic example shows that
exceptional initial points and tangential collapse cannot simply be omitted.
\end{abstract}

\section{Source question and answer}

The source is Hinz, T\"olle, and Viitasaari~\cite{HTV}.  In Section 1.2,
PDF page 6, after their Doss-transform existence theory, they write that the
main open problem is to prove, under reasonable assumptions, that ``the
integral process itself will be variable,'' with the goal of using standard
fixed-point arguments.

\begin{figure}[ht]
  \centering
  \includegraphics[width=\textwidth]{figures/open_problem_crop.png}
  \caption{The source statement on PDF page 6.  This is a direct crop of the
  source paper.}
\end{figure}

This packet now gives two complementary advances.  First, bounded compression
forces variability for almost every starting point with no geometric
regularity of the singular set.  This applies to the full multidimensional
Doss class of the source and, in a different time-regularity regime, to
regular Lagrangian flows without Doss structure.  Second, for one
codimension-one jump interface it gives pointwise fixed-point constructions.
The rough interface theorem keeps $Y\in C^\gamma$, $\gamma>1/2$, while the
absolutely continuous theorem allows a state-dependent normal velocity.

\section{Definitions and setup}

Let $H=\{x\in\R^n:h(x)=0\}$, where
\[
 h(x)=\nu\cdot x-a,\qquad |\nu|=1.
\]
Given two matrix fields $\sigma_\pm:\R^n\to\R^{n\times n}$, set
\begin{equation}\label{eq:sigma-piecewise}
 \sigma(x)=
 \begin{cases}
  \sigma_+(x),&h(x)\geq0,\\
  \sigma_-(x),&h(x)<0.
 \end{cases}
\end{equation}
The representative chosen on $H$ is irrelevant to all distributional
derivatives and to the integrals below, because a solution meets $H$ at most
once.

For a scalar $\varphi\in\BV(\R^n)$, a continuous path
$X:[0,T]\to\R^n$ is $(s,p)$-variable with respect to $\varphi$ if some
relatively compact open neighborhood $U$ of $X([0,T])$ satisfies
\[
 t\longmapsto \int_U |X_t-z|^{-n+1-s}\,\lVert D\varphi\rVert(dz)
 \quad\text{in }L^p(0,T).
\]
For a matrix coefficient this is required componentwise, and we write
$X\in\V(\sigma,s,p)$.  The source notation $\V(\sigma,s)$ means $p=1$.

We study
\begin{equation}\label{eq:main-equation}
 X_t=x_0+\int_0^t\sigma(X_u)\,dY_u.
\end{equation}
For H\"older drivers, integrals on intervals that do not contain the crossing
are ordinary Young integrals.  The global integral is their additive
concatenation.  It equals the source's generalized Lebesgue--Stieltjes
integral: the integrand is H\"older on each side, has at most one jump in
time, changing its value at the crossing does not alter either integral, and
the restriction property joins the two pieces.

\section{Idea of the proof}

For bounded compression, the proof is a one-line interchange of integrals
with a decisive consequence.  If $\Phi_t$ pushes Lebesgue measure forward
with density at most $L$, then for $\mu=|D\varphi|\restriction_U$,
\[
 \int_B\!\int_0^T\!\int_U
 |\Phi_t(x)-z|^{-n+1-s}\,\mu(dz)\,dt\,dx
 \leq L T\!\int_U\!\int_K |y-z|^{-n+1-s}\,dy\,\mu(dz).
\]
The last integral is finite because $n-1+s<n$.  Hence the variability energy
is finite for almost every initial point.  Taking a countable sequence
$s_j\uparrow1$ gives one exceptional null set valid for every $s<1$.

The source's Doss trajectories are conjugates of translations by a fixed
bi-Lipschitz map, hence have bounded compression.  Independently, Ambrosio's
regular Lagrangian flows have bounded compression by definition.  These two
facts turn the abstract estimate into the generic-start results below.

The decisive quantity is the signed distance $h(X_t)$.  A common-row identity
for $\sigma_+$ and $\sigma_-$ makes this coordinate independent of the unknown
path:
\[
 \nu^\top\sigma_\pm(x)=\lambda^\top
 \quad\Longrightarrow\quad
 h(X_t)=h(x_0)+\lambda\cdot(Y_t-Y_0).
\]
If the right-hand side is a strictly increasing bi-Lipschitz clock, the side
of the interface occupied at every time is known before solving the equation.
One therefore solves an ordinary smooth Young equation with $\sigma_-$ before
the unique crossing and one with $\sigma_+$ afterward.  Standard Picard
contractions give the two pieces and uniqueness.

The same identity proves variability.  Locally, the derivative measure of a
piecewise smooth coefficient is bounded by volume measure plus surface measure
on $H$.  The volume contribution to the relevant Riesz potential is uniformly
bounded.  The surface contribution at a point $x$ has size
\[
 \int_H |x-z|^{-n+1-s}\,d\HH^{n-1}(z)
 \asymp \dist(x,H)^{-s}.
\]
Linear crossing turns this into $|t-\tau|^{-s}$ near the unique hitting time
$\tau$, and its $p$th power is integrable exactly when $sp<1$.

For an absolutely continuous driver the normal velocity need not be a fixed
row.  It is enough that both one-sided vector fields have uniformly positive
normal speed.  Candidate paths with that same speed form a compact convex set.
Although $\sigma$ is discontinuous, uniform convergence of candidate paths
gives pointwise convergence of the integrands away from the single crossing
time.  Dominated convergence makes the Picard operator continuous, so
Schauder applies.

\section{The bounded-compression principle}

The next theorem has no interface or rectifiability hypothesis.  It is the
strongest result in the packet.

\begin{theorem}[Bounded compression forces generic variability]
\label{thm:compression}
Let $B\subset\R^n$ be bounded and measurable and let
$\Phi:[0,T]\times B\to\R^n$ be measurable, with each trajectory
$t\mapsto\Phi_t(x)$ continuous.  Suppose that
\begin{equation}\label{eq:compression}
 (\Phi_t)_\#(\mathbf1_B\mathcal L^n)\leq L\mathcal L^n
 \qquad(0\leq t\leq T)
\end{equation}
for a constant $L$, and that all the trajectories take values in one bounded
set $K$.  Let $\varphi\in BV_{\rm loc}(\R^n)$.  Then, outside one
Lebesgue-null subset of $B$, the path $t\mapsto\Phi_t(x)$ is
$(s,1)$-variable with respect to $\varphi$ simultaneously for every
$s\in(0,1)$.  The same conclusion holds componentwise for any finite matrix
of locally $BV$ functions.
\end{theorem}

\begin{proof}
Choose a bounded open set $U$ containing the closure of $K$ and put
$\mu=|D\varphi|\restriction_U$.  For fixed $s\in(0,1)$ define
\[
 E_s(x)=\int_0^T\int_U
 |\Phi_t(x)-z|^{-n+1-s}\,\mu(dz)\,dt.
\]
Tonelli's theorem and~\eqref{eq:compression} give
\begin{align}
 \int_B E_s(x)\,dx
 &\leq L\int_0^T\int_U\int_K
 |y-z|^{-n+1-s}\,dy\,\mu(dz)\,dt.\label{eq:energy-average}
\end{align}
If $K\cup U$ has diameter at most $R$, translation and polar coordinates
give, uniformly in $z\in U$,
\[
 \int_K|y-z|^{-n+1-s}\,dy
 \leq |\mathbb S^{n-1}|\int_0^R r^{-s}\,dr
 =\frac{|\mathbb S^{n-1}|R^{1-s}}{1-s}.
\]
Since $\mu(U)<\infty$, the right-hand side of
\eqref{eq:energy-average} is finite.  Thus $E_s(x)<\infty$ for almost every
$x$, which is precisely $(s,1)$-variability.

To obtain a single null set, apply the preceding argument to a countable
sequence $s_j\uparrow1$.  Given $s<s_j$, split distances at one; on the
bounded set $U$ the kernel of exponent $n-1+s$ is bounded by a constant times
one plus the kernel of exponent $n-1+s_j$.  Hence finiteness for $s_j$
implies finiteness for $s$.  A finite intersection handles all entries of a
matrix coefficient.
\end{proof}

\begin{corollary}[Generic starts in the full Doss class]
\label{cor:doss-generic}
Assume the multidimensional hypotheses of the source that produce a
bi-Lipschitz Doss map $f:\R^n\to\R^n$ satisfying
$\nabla f=\sigma\circ f$: namely its Assumptions 3.12 and 3.15 together with
its angle condition (3.10).  Let $Y\in C^\gamma([0,T];\R^n)$ with
$\gamma>1/2$.  Then for Lebesgue-almost every initial point $x\in\R^n$,
\begin{equation}\label{eq:doss-flow}
 X_t^x=f\bigl(Y_t-Y_0+f^{-1}(x)\bigr)
\end{equation}
belongs to $\V(\sigma,s)$ for every $s<1$.  Consequently it is the source's
variability solution, and hence its unique variability solution under the
source's uniqueness hypotheses.  No upper-regularity, occupation-density, or
Riesz-energy assumption on the driver is required.
\end{corollary}

\begin{proof}
Set $\Phi_t=f\circ\tau_{Y_t-Y_0}\circ f^{-1}$.  Its inverse is
$f\circ\tau_{Y_0-Y_t}\circ f^{-1}$, so
\[
 \operatorname{Lip}(\Phi_t^{-1})
 \leq\operatorname{Lip}(f)\operatorname{Lip}(f^{-1})
\]
uniformly in $t$.  The standard measure estimate for a Lipschitz map gives
\eqref{eq:compression} on every bounded set $B$, with
$L\leq(\operatorname{Lip}(f)\operatorname{Lip}(f^{-1}))^n$.  The range of
$[0,T]\times B$ is bounded because $Y([0,T])$ and $B$ are bounded and $f$
is bi-Lipschitz.  Theorem~\ref{thm:compression}, first on $B(0,m)$ and then
for all $m\in\mathbb N$, proves the simultaneous variability assertion.

This argument is not circular: formula~\eqref{eq:doss-flow} defines a
continuous path before an integral is invoked.  Bounded compression proves
its variability.  One may then follow the second half of the source's
Lemma 5.7: the $BV$ composition estimate for the bi-Lipschitz inverse transfers
variability of $X$ with respect to $\sigma$ to variability of
$Y-Y_0+f^{-1}(x)$ with respect to each entry of $\nabla f=\sigma\circ f$.
The source's Theorem 2.13 then applies the change-of-variables formula and
identifies $X$ with a solution of~\eqref{eq:main-equation}.  The path is
$\gamma$-H\"older.  Choose $(1-\gamma)/\gamma<s<1$ to meet Definition 3.1
of the source.  Its uniqueness theorem then applies whenever its stated Doss
hypotheses are in force.
\end{proof}

The same averaging lemma also combines with a completely different existence
theory and thereby removes the Doss assumptions, at the price of a smoother
driver and an almost-everywhere flow notion.

\begin{corollary}[A non-Doss regular-Lagrangian-flow theorem]
\label{cor:ambrosio}
Let $\sigma:\R^n\to\R^{n\times n}$ be bounded with locally $BV$ entries and
write $\sigma^{(k)}$ for its $k$th column.  Assume
\[
 \operatorname{div}\sigma^{(k)}\in L^\infty(\R^n),
 \qquad k=1,\ldots,n.
\]
For $Y\in W^{1,\infty}([0,T];\R^n)$, the vector field
$b_t(x)=\sigma(x)\dot Y_t$ has a unique regular Lagrangian flow $X_t(x)$.
For almost every $x$, its trajectory belongs to
$\V(\sigma,s)$ simultaneously for every $s<1$ and solves
\[
 X_t(x)=x+\int_0^t\sigma(X_u(x))\dot Y_u\,du.
\]
In particular, regarding $Y$ as $C^\gamma$ for any
$\gamma\in(1/2,1)$ and choosing $s\in(1-\gamma,1)$, this trajectory is a
variability solution in the sense of the source.
\end{corollary}

\begin{proof}
The field $b$ belongs to $L^1_tBV_{\rm loc,x}$, is bounded, and
\[
 \operatorname{div}b_t
 =\sum_{k=1}^n\dot Y_t^k\operatorname{div}\sigma^{(k)}
 \in L^1(0,T;L^\infty(\R^n)).
\]
Its boundedness supplies Ambrosio's growth condition.  The regular
Lagrangian flow theorem~\cite{Ambrosio,AmbrosioOverview} therefore gives existence,
uniqueness in the regular-flow class, and a uniform compression bound.  For
almost every initial point its trajectory is absolutely continuous; discard
the exceptional null set, which can only decrease each pushforward measure.
For
initial points in a bounded $B$, the estimate
$|X_t(x)-x|\leq\|\sigma\|_\infty\|\dot Y\|_{L^1}$ places the trajectories
in a common bounded set.  Theorem~\ref{thm:compression} applies to every
matrix entry.  Finally the trajectories are Lipschitz, so in Definition 3.1
one may take path exponent $\alpha=1$; the required inequality is then
$s>1-\gamma$.
\end{proof}

\section{The rough-driver fixed-point theorem}

We first recall the standard smooth Young equation in the precise strength
needed here.

\begin{lemma}[Young--Picard lemma]\label{lem:young-picard}
Let $\gamma\in(1/2,1)$, $Y\in C^\gamma([r,q];\R^n)$, and
$V\in C_b^2(\R^n;\R^{n\times n})$.  For every $z\in\R^n$ there is a unique
$Z\in C^\gamma([r,q];\R^n)$ satisfying
\[
 Z_t=z+\int_r^t V(Z_u)\,dY_u.
\]
It is obtained by Picard contraction on sufficiently short subintervals and
concatenation.  If $V$ is bounded, the concatenation reaches all of $[r,q]$.
\end{lemma}

\begin{proof}
For $f\in C^\gamma$ the Young--Loeve estimate gives, on an interval $I$,
\[
 \left|\int_u^v f\,dY-f_u(Y_v-Y_u)\right|
 \leq C_\gamma[f]_\gamma[Y]_\gamma|v-u|^{2\gamma}.
\]
On a $C^\gamma$ ball of paths with a common initial value, boundedness of
$DV$ and $D^2V$ gives
\begin{align*}
 \|V(X)-V(Z)\|_\infty&\leq \|DV\|_\infty\|X-Z\|_\infty,\\
 [V(X)-V(Z)]_\gamma
 &\leq C_V\bigl([X-Z]_\gamma+
 ([X]_\gamma+[Z]_\gamma)\|X-Z\|_\infty\bigr).
\end{align*}
Since $X$ and $Z$ agree at the left endpoint,
$\|X-Z\|_\infty\leq |I|^\gamma[X-Z]_\gamma$.  The Young--Loeve estimate
therefore makes the integral map a contraction when $|I|$ is small enough;
the same estimate makes the chosen ball invariant.  Iteration constructs the
local solution.  Uniqueness allows adjacent pieces to concatenate.  The
bounded vector field supplies a uniform local range bound, so finitely many
pieces cover the compact time interval.
\end{proof}

\begin{theorem}[Rough clock-row theorem]\label{thm:rough}
Let $\gamma\in(1/2,1)$ and $Y\in C^\gamma([0,T];\R^n)$.  Suppose
$\sigma_\pm\in C_b^2(\R^n;\R^{n\times n})$ and there is
$\lambda\in\R^n$ such that
\begin{equation}\label{eq:common-row}
 \nu^\top\sigma_+(x)=\nu^\top\sigma_-(x)=\lambda^\top
 \qquad(x\in\R^n).
\end{equation}
Assume the scalar control
\begin{equation}\label{eq:g-def}
 g(t)=h(x_0)+\lambda\cdot(Y_t-Y_0)
\end{equation}
is increasing and bi-Lipschitz: for some $0<m\leq M$,
\begin{equation}\label{eq:bilip-clock}
 m(t-u)\leq g(t)-g(u)\leq M(t-u),\qquad 0\leq u<t\leq T.
\end{equation}
Then equation~\eqref{eq:main-equation} has a unique
$X\in C^\gamma([0,T];\R^n)$ solution, constructed entirely by standard
Young--Picard fixed points.  Moreover,
\begin{equation}\label{eq:sp-conclusion}
 X\in\V(\sigma,s,p)
 \quad\text{for every }s\in(0,1),\ p\in[1,\infty)
 \text{ with }sp<1.
\end{equation}
In particular, choosing
\[
 \frac{1-\gamma}{\gamma}<s<1
\]
makes $X$ a variability solution for $\sigma$ and $Y$ in the exact sense of
the source.
\end{theorem}

\begin{proof}
The strictly increasing function $g$ has at most one zero.  If it has a zero,
call it $\tau$; otherwise the sign is constant.  On every interval on which
$g<0$, solve the smooth Young equation with coefficient $\sigma_-$ using
Lemma~\ref{lem:young-picard}; on every interval on which $g>0$, use
$\sigma_+$.  If there is a crossing, use the endpoint of the first solution as
the initial value of the second.

Taking the scalar product with $\nu$ in either smooth equation and using
\eqref{eq:common-row} gives
\[
 h(X_t)=h(x_0)+\int_0^t\lambda\cdot dY_u=g(t).
\]
Thus the constructed path lies in the half-space whose coefficient was used,
so its concatenation solves the piecewise equation.  The possible discrepancy
at $t=\tau$ is a one-point change of the integrand and has zero integral.

Conversely, any solution of the piecewise equation has the same normal
coordinate $g$.  Its switching schedule is therefore identical.  On each
side it solves the corresponding smooth Young equation, where
Lemma~\ref{lem:young-picard} gives uniqueness.  This proves global uniqueness.

It remains to prove~\eqref{eq:sp-conclusion}.  Let $U\Subset\R^n$ be a bounded
open neighborhood of $X([0,T])$.  The $BV$ product rule for
\eqref{eq:sigma-piecewise} and the bounded first derivatives of the two
branches give, for every scalar entry $\sigma_{jk}$,
\begin{equation}\label{eq:measure-majorization}
 |D\sigma_{jk}|\!\restriction_U
 \leq C_U\,\mathcal L^n\!\restriction_U
   +J_U\,\HH^{n-1}\!\restriction_{H\cap U}.
\end{equation}
The volume part has uniformly bounded potential because $s<1$:
\[
 \sup_{x\in X([0,T])}\int_U|x-z|^{-n+1-s}\,dz<\infty.
\]
For $x\notin H$, rotate and translate coordinates, enlarge $H\cap U$ to the
whole hyperplane, and scale the tangential variable to obtain
\begin{align}\label{eq:hyperplane-potential}
 \int_{H\cap U}|x-z|^{-n+1-s}\,d\HH^{n-1}(z)
 &\leq \int_{\R^{n-1}}
 (|w|^2+|h(x)|^2)^{-(n-1+s)/2}\,dw\\
 &=C_{n,s}|h(x)|^{-s}.
\end{align}
For $n=1$ this is the same identity with $\HH^0$ on a point.

If $g$ never vanishes, compactness gives $|g|\geq\delta>0$.  If it vanishes
at $\tau$,~\eqref{eq:bilip-clock} gives
\[
 |h(X_t)|=|g(t)|\geq m|t-\tau|.
\]
Combining this with~\eqref{eq:measure-majorization} and
\eqref{eq:hyperplane-potential} bounds each variability potential by
$C(1+|t-\tau|^{-s})$.  Its $p$th power is integrable precisely when $sp<1$;
the infinite value at $t=\tau$ is immaterial.  This proves
\eqref{eq:sp-conclusion}.

Finally, $X$ is $\gamma$-H\"older.  Because $\gamma>1/2$, the interval
$((1-\gamma)/\gamma,1)$ is nonempty.  For any $s$ in it we have
$X\in\V(\sigma,s)$ and $\gamma s+\gamma>1$, exactly the source definition of
a variability solution.  On each side the generalized integral agrees with
the Young integral, and additivity across the single crossing proves that the
constructed equation is the source equation.
\end{proof}

\begin{remark}[A concrete rough family]
After a linear change of coordinates, take $\nu=\lambda=e_1$ and
$Y_t=(t,W_t)$, where $W$ is any $(n-1)$-dimensional $\gamma$-H\"older path,
$\gamma>1/2$.  Any two bounded $C^2$ branches whose first matrix row is
$e_1^\top$ satisfy the theorem, regardless of the size of their jump in the
remaining rows.  Thus the theorem is not merely an ordinary smooth-time ODE:
all tangential equations may be driven by a genuinely rough signal.
\end{remark}

\section{A direct Schauder theorem for absolutely continuous drivers}

The next theorem allows the normal row to depend on the point and to differ on
the two sides.  A common positive normal velocity replaces the algebraic
identity~\eqref{eq:common-row}.

\begin{theorem}[Common-crossing Schauder theorem]\label{thm:schauder}
Let $Y\in W^{1,\infty}([0,T];\R^n)$ and let the bounded fields
$\sigma_\pm:\R^n\to\R^{n\times n}$ be locally Lipschitz.  Assume there are
constants $0<c\leq M<\infty$ such that, for almost every $t$ and every
$x\in\R^n$,
\begin{equation}\label{eq:crossing}
 |\sigma_\pm(x)\dot Y_t|\leq M,
 \qquad
 \nu\cdot\sigma_\pm(x)\dot Y_t\geq c.
\end{equation}
Then the discontinuous integral equation
\[
 X_t=x_0+\int_0^t\sigma(X_u)\dot Y_u\,du
\]
has a unique absolutely continuous solution.  It is obtained as a fixed point
of the global discontinuous Picard map by the ordinary Schauder theorem, and
it satisfies
\[
 X\in\V(\sigma,s,p)
 \qquad(0<s<1,\ 1\leq p<\infty,\ sp<1).
\]
\end{theorem}

\begin{proof}
Work in $C([0,T];\R^n)$ and define $K$ as the set of paths $Z$ such that
$Z_0=x_0$ and, for $0\leq u<t\leq T$,
\begin{equation}\label{eq:K}
 |Z_t-Z_u|\leq M(t-u),
 \qquad h(Z_t)-h(Z_u)\geq c(t-u).
\end{equation}
The set is nonempty (it contains $x_0+ct\nu$), convex, closed, uniformly
bounded, and equicontinuous.  Hence it is compact by Arzel\`a--Ascoli.

Define
\[
 (\Tmap Z)_t=x_0+\int_0^t\sigma(Z_u)\dot Y_u\,du.
\]
The two inequalities in~\eqref{eq:crossing} show directly that
$\Tmap(K)\subset K$.

We claim that $\Tmap:K\to K$ is continuous even though $\sigma$ is not.  If
$Z^k\to Z$ uniformly in $K$, then $t\mapsto h(Z_t)$ is strictly increasing by
\eqref{eq:K}; consequently $Z$ meets $H$ at at most one time.  At every other
time, $Z_t$ lies strictly in one half-space, $Z^k_t$ eventually lies in the
same half-space, and local continuity of that branch gives
\[
 \sigma(Z^k_t)\dot Y_t\longrightarrow\sigma(Z_t)\dot Y_t
\]
for almost every $t$.  The difference of the integrands is bounded by $2M$.
Dominated
convergence therefore yields
\[
 \|\Tmap Z^k-\Tmap Z\|_\infty
 \leq\int_0^T
 |\sigma(Z^k_t)\dot Y_t-\sigma(Z_t)\dot Y_t|\,dt\longrightarrow0.
\]
Schauder's theorem supplies a fixed point $X\in K$.

For uniqueness, every absolutely continuous solution automatically belongs to
$K$.  On each open time interval on which it stays in one half-space, it solves
a time-dependent locally Lipschitz ODE.  Two solutions agree by Gronwall until
the first one reaches $H$; since their paths then agree, their hitting times
and hitting points agree.  After the crossing they solve the same one-sided
ODE from the same initial value.  Thus they coincide globally.

Finally,~\eqref{eq:K} implies that $h(X_t)$ crosses zero at most once and, at a
crossing time $\tau$, $|h(X_t)|\geq c|t-\tau|$.  The measure estimate
\eqref{eq:measure-majorization} remains valid for locally Lipschitz branches on
the bounded neighborhood $U$, and the argument
\eqref{eq:hyperplane-potential} gives the stated $(s,p)$-variability.
\end{proof}

\section{Why transversality is doing real work}

The conclusion is false if one keeps strong matrix nondegeneracy but drops
output crossing.  The companion packet proves the following exact example;
we repeat its algebraic core to delimit Theorems~\ref{thm:rough} and
\ref{thm:schauder}.

\begin{proposition}[Tangential-collapse obstruction]\label{prop:obstruction}
In $\R^2$ let $H=\{x_2=0\}$ and
\[
 A_+=\begin{pmatrix}2&-1\\-1&1\end{pmatrix},\qquad
 A_-=\begin{pmatrix}\frac32&-\frac12\\-\frac12&\frac12\end{pmatrix},
 \qquad
 \sigma=A_+\mathbf1_{\{x_2\geq0\}}+A_-\mathbf1_{\{x_2<0\}}.
\]
The coefficient is bounded, symmetric, and uniformly positive definite; its
inverse has distributionally curl-free rows, and it satisfies the determinant
and angle assumptions in the source's multidimensional Doss theory.  For every
$s\in(0,1)$ and finite $p$, choose
$0<\alpha<\min\{1,(sp)^{-1}\}$ and set
\[
 X_t=(0,t^\alpha),\qquad Y_t=(t,t).
\]
Then $X\in\V(\sigma,s,p)$, but
\[
 \int_0^t\sigma(X_u)\,dY_u=(t,0)
\]
lies on $H$ and is not $(\rho,q)$-variable for any $\rho>0$ and any
$q\geq1$.
\end{proposition}

\begin{proof}
The eigenvalues of $A_+$ are $(3\pm\sqrt5)/2$ and those of $A_-$ are
$1\pm1/\sqrt2$.  Both matrices are therefore uniformly positive definite.
Their inverses are
\[
 A_+^{-1}=\begin{pmatrix}1&1\\1&2\end{pmatrix},\qquad
 A_-^{-1}=\begin{pmatrix}1&1\\1&3\end{pmatrix};
\]
because they differ only in the $(2,2)$ entry, the inverse rows have zero
distributional curl across the horizontal line.  Also
$A_+(1,1)^\top=A_-(1,1)^\top=(1,0)^\top$.

The jump measure is a nonzero multiple of $\HH^1\restriction_H$.  At distance
$t^\alpha$ from $H$, its order-$(1-s)$ potential is $Ct^{-\alpha s}$, which is
in $L^p$ when $\alpha sp<1$.  The Picard image is $(t,0)$ by the displayed
matrix identity.  On $H$ the potential contains
$\int_{-\varepsilon}^{\varepsilon}|r|^{-1-\rho}dr=\infty$ at every time,
so the image is in no positive variability class.
\end{proof}

\begin{corollary}[The generic-start qualifier is sharp]
\label{cor:exceptional}
For the coefficient in Proposition~\ref{prop:obstruction}, let
$B\in C^\gamma([0,T])$ be nonconstant with $\gamma>1/2$, put
$Y=(B,B)$, and start at any $x\in H$.  Then
\[
 X_t=x+(B_t-B_0,0)
\]
is a solution that remains in $H$ and is not $(s,p)$-variable for any
$s>0$ and $p\geq1$.  Moreover, there is no variability solution from such
an initial point.  Thus Theorem~\ref{thm:compression} cannot in general be
improved from almost every start to every start, even for a uniformly elliptic
coefficient satisfying the source's Doss structural conditions.
\end{corollary}

\begin{proof}
Both matrices satisfy
$A_+(1,1)^\top=A_-(1,1)^\top=e_1$.  On the displayed candidate the chosen
representative is the constant matrix $A_+$, so the Young integral is exactly
$e_1(B_t-B_0)$ and the equation holds.  Its range lies in $H$, where the
Riesz potential of the nonzero jump measure is infinite at every time.

If a variability solution existed, all its component integrals would exist
and their linearity, together with the same column-sum identity on both sides,
would force it to equal the displayed candidate.  That candidate is not
variable, a contradiction.
\end{proof}

Thus uniform ellipticity, smooth driving, and input variability do not make the
Picard map preserve variability.  The crossing hypotheses in the two main
theorems prevent exactly the tangential collapse exhibited above.

\section{Sharpness, limitations, and novelty}

Theorem~\ref{thm:compression} is optimal in two immediate senses.  The
restriction $s<1$ is exactly the local-integrability threshold for its Riesz
kernel, and Corollary~\ref{cor:exceptional} exhibits an entire null
hyperplane of initial points whose trajectories fail every positive
variability condition.  Bounded compression controls spatial averages; it
cannot prohibit a singular exceptional orbit.

The threshold $sp<1$ is sharp for a genuine jump and a linear crossing.  If a
component has nonzero jump trace at the hitting point, the hyperplane potential
is comparable from below to $|t-\tau|^{-s}$, which is not in $L^p$ when
$sp\geq1$.  The $p=1$ range nevertheless contains every $s<1$, and hence the
source compatibility interval $s>(1-\gamma)/\gamma$ for every
$\gamma>1/2$.

The rough theorem deliberately assumes a single flat interface, smooth
one-sided branches, a common normal row, and a bi-Lipschitz clock projection.
The tangential driver can be fully H\"older-rough, but a driver with no monotone
projection is not covered.  Curved or multiple interfaces should be approachable
with event functions and finitely many transversal hits, but their crossing
schedule is no longer algebraically fixed and is not claimed here.  The
Schauder theorem allows state-dependent normal motion but requires an
absolutely continuous driver.

The generic Doss corollary is broad in the coefficient and driver but still
uses the Doss transform to construct the trajectory.  The Ambrosio corollary
does not use Doss and permits noninvertible coefficients, but it assumes a
Lipschitz driver, bounded column divergences, and obtains a regular
Lagrangian flow only for almost every initial point.  Therefore neither
corollary is the requested general rough-driver fixed-point theorem at an
arbitrary prescribed start.  They instead identify bounded compression as a
powerful sufficient mechanism for the authors' central output-variability
property.

A bounded search on 26 August 2026 covered the run's lightweight indexes, the
exact source phrase and arXiv id, the authors' follow-ups, and combinations of
``variability,'' ``fixed point,'' ``bounded compression,'' ``regular
Lagrangian flow,'' ``hyperplane interface,'' and ``transversality.''  No
matching bounded-compression-to-variability theorem or fixed-point-plus-
variability theorem was found.  Ambrosio~\cite{Ambrosio} supplies the regular
Lagrangian flow theorem used in Corollary~\ref{cor:ambrosio}, but not the
Riesz-energy averaging conclusion.  Burden--Sastry--Koditschek--Revzen~\cite{Burden} prove well-defined
piecewise-differentiable flows for event-selected discontinuous vector fields
under transversality, but do not formulate the Riesz-potential variability
conclusion or the rough clock-row theorem.  L\'opez Pouso~\cite{Pouso} uses
Schauder ideas for discontinuous boundary-value problems; its null-contact-set
continuity mechanism is philosophically close, but the equations and
variability conclusion are different.  The source itself proves that smooth
curves with finitely many transversal hits are variable, but it does not use
that fact to construct a fixed point.  The present contribution is the
combination: fixed-point construction, forced crossing geometry, and the sharp
$sp<1$ estimate.

Because the search is bounded rather than a systematic citation review, the
proper label is \emph{candidate substantial partial result, likely valid}.
The bounded-compression result covers the source's full Doss coefficient class
only generically in the initial point, and the pointwise fixed-point results
remain restricted to one interface.  This is not a full answer to the source
problem.

\section{Verification checklist}

The proof reduces to seven independently checkable mechanisms:
\begin{enumerate}
\item Tonelli plus bounded compression changes the initial-point average of
the variability energy into an ordinary locally integrable Riesz kernel.
\item Bi-Lipschitz conjugates of translations have a uniform compression
constant, so the estimate applies to the source's Doss flow before its chain
rule is used.
\item Ambrosio's theorem applies to
$b_t=\sum_k\dot Y_t^k\sigma^{(k)}$ under the stated bounded-divergence and
growth hypotheses.
\item The common-row identity forces the normal coordinate of every solution
to equal the prescribed clock $g$.
\item The fixed crossing schedule reduces the rough discontinuous equation to
at most two standard smooth Young equations, each obtained by Picard
contraction.
\item Surface measure on a hyperplane has Riesz potential exactly of order
$\dist(\cdot,H)^{-s}$; linear crossing gives the sharp condition $sp<1$.
\item In the absolutely continuous case, a candidate path contacts the jump
set at most once, which is enough for dominated convergence and continuity of
the global Picard map.
\end{enumerate}
No numerical evidence is used as proof.

\begin{thebibliography}{9}

\bibitem{HTV}
M. Hinz, J. M. T\"olle, and L. Viitasaari,
\emph{Variability of paths and differential equations with
$BV$-coefficients}, Ann. Inst. H. Poincar\'e Probab. Statist. 59 (2023),
2036--2082; \href{https://arxiv.org/abs/2003.11698}{arXiv:2003.11698}.

\bibitem{Ambrosio}
L. Ambrosio,
\emph{Transport equation and Cauchy problem for $BV$ vector fields},
Invent. Math. 158 (2004), 227--260;
\href{https://doi.org/10.1007/s00222-004-0367-2}{doi:10.1007/s00222-004-0367-2}.

\bibitem{AmbrosioOverview}
L. Ambrosio,
\emph{Transport equation and Cauchy problem for $BV$ vector fields and
applications}, Journ\'ees \'equations aux d\'eriv\'ees partielles (2004),
Expos\'e no.~1;
\href{https://doi.org/10.5802/jedp.1}{doi:10.5802/jedp.1}.

\bibitem{Burden}
S. A. Burden, S. S. Sastry, D. E. Koditschek, and S. Revzen,
\emph{Event-selected vector field discontinuities yield
piecewise-differentiable flows},
\href{https://arxiv.org/abs/1407.1775}{arXiv:1407.1775}.

\bibitem{Pouso}
R. L\'opez Pouso,
\emph{Schauder's fixed-point theorem: new applications and a new version for
discontinuous operators}, Boundary Value Problems 2012, Article 92;
\href{https://doi.org/10.1186/1687-2770-2012-92}{doi:10.1186/1687-2770-2012-92}.

\end{thebibliography}

\end{document}
